\section{Level A: bridging}\label{sec:levela}

\subsection{Model and estimator}

Let $\Omega\subseteq U\times J$ be the observed (reviewer, item) pairs of an epoch, with
$|U|=n$ reviewers and $|J|=m$ items, and let $r_{uj}\in[0,1]$ be the probability that
reviewer $u$ declares for item $j$. Each reviewer carries a weight $w_u\ge0$ computed from
the previous epoch: $0$ during probation, $1$ for the founders who bootstrap the network, and
$\min(w_{\max},E_u)$ afterwards, where $E_u$ is the evaluator score of \cref{sec:levelc}. The design then applies the anti-collusion discount of
\cref{sec:collusion}; the engine implements it, but the protocol does not yet apply it when it
computes these weights. The model is the
one-dimensional ($d=1$) factorization
\begin{equation}\label{eq:mf}
  \hat r_{uj} \;=\; \mu + \iota_u + \iota_j + f_u f_j ,
\end{equation}
with a global mean $\mu$, reviewer and item intercepts $\iota_u,\iota_j$ (written
$b_u,b_j$ in the repository), a reviewer position $f_u$ on a latent axis, and an item loading
$f_j$. The parameters $\xi=(\mu,\iota_U,\iota_J,f_U,f_J)$ minimize
\begin{equation}\label{eq:loss}
  L(\xi) = \sum_{(u,j)\in\Omega} w_u\,(r_{uj}-\hat r_{uj})^2
  + \lambda_b\Bigl(\sum_u \iota_u^2+\sum_j\iota_j^2\Bigr)
  + \lambda_f\Bigl(\sum_u f_u^2+\sum_j f_j^2\Bigr),
\end{equation}
with $\lambda_b=0.15$ and $\lambda_f=0.03$, and $\mu$ unpenalized. The objective is
non-convex. The reference implementation minimizes it with L-BFGS and a strong-Wolfe line
search~\cite{liu1989lbfgs,nocedal2006numerical} from eight seeded starts, keeps the lowest
objective, and fixes the sign of the axis by making the largest $|f_j|$ non-negative.

The \emph{bridge score} is pessimistic: $B_j=\min_s \iota_j^{(s)}$, taken over the full fit and
$10$ refits on subsamples that keep each observation with probability $0.85$~\cite{efron1993bootstrap}.
The gate compares $B_j$ with a threshold $\tau\approx0.08$ and a band half-width
$\varepsilon\approx0.008$, both provisional (\cref{sec:setting}).

The design rationale, taken from Community Notes~\cite{wojcik2022birdwatch}, is that heavy
penalties on the intercepts force the model to explain approval through the factor term first.
Approval that a reviewer's side alone explains goes into $f_uf_j$. Only approval that no side
explains, that is, approval \emph{across} the divide, survives in $\iota_j$. The rest of this
section examines what the objective actually guarantees.

\subsection{Invariances of the objective}

The data term of \eqref{eq:loss} depends on $\xi$ only through the predictions
$\hat r_{uj}$. Several transformations leave every prediction unchanged:
\begin{align*}
  &\text{(G1)}\;\; \mu\mapsto\mu+c,\;\; \iota_j\mapsto\iota_j-c\;\;\forall j;
  &&\text{(G2)}\;\; \mu\mapsto\mu+c,\;\; \iota_u\mapsto\iota_u-c\;\;\forall u;\\
  &\text{(G3)}\;\; f_u\mapsto f_u+c\;\;\forall u,\;\; \iota_j\mapsto\iota_j-c f_j\;\;\forall j;
  &&\text{(G4)}\;\; f_j\mapsto f_j+c\;\;\forall j,\;\; \iota_u\mapsto\iota_u-c f_u\;\;\forall u;\\
  &\text{(G5)}\;\; f_u\mapsto s f_u,\;\; f_j\mapsto f_j/s,\;\; s\neq0.
\end{align*}
(G3) translates the reviewers' axis and (G4) the items' axis; (G5) rescales the axis and, for
$s=-1$, flips its sign. Along these directions only the penalties vary, so the penalties alone
decide where the solution sits. This yields exact identities.

\begin{lemma}[Stationarity identities]\label{lem:gauge}\proved{}
For any weights $w_u\ge0$ and any stationary point of \eqref{eq:loss} (in particular any local
minimum),
\begin{align}
  &\textstyle\sum_j \iota_j = 0 \quad\text{and}\quad \sum_u \iota_u = 0, \label{eq:zerosum}\\
  &\textstyle\lambda_b\sum_j f_j\,\iota_j \;=\; \lambda_f\sum_u f_u, \label{eq:gauge-rev}\\
  &\textstyle\lambda_b\sum_u f_u\,\iota_u \;=\; \lambda_f\sum_j f_j, \label{eq:gauge-item}\\
  &\textstyle\sum_u f_u^2 \;=\; \sum_j f_j^2. \label{eq:scale}
\end{align}
\end{lemma}
\begin{proof}
Let $\xi^*$ be stationary and $\xi(c)$ a curve through it along one of the families above. The
data term is constant along the curve, so $0=\frac{d}{dc}L(\xi(c))|_{c=0}$ is the derivative of
the penalty alone. For (G1) this derivative is $-2\lambda_b\sum_j\iota_j$, and (G2) is
symmetric. For (G3) it is
$\frac{d}{dc}\bigl[\lambda_b\sum_j(\iota_j-cf_j)^2+\lambda_f\sum_u(f_u+c)^2\bigr]_{c=0}
=-2\lambda_b\sum_j f_j\iota_j+2\lambda_f\sum_u f_u$, and (G4) is symmetric. For (G5) it is
$\frac{d}{ds}\bigl[\lambda_f(s^2\sum_u f_u^2+s^{-2}\sum_j f_j^2)\bigr]_{s=1}
=2\lambda_f(\sum_u f_u^2-\sum_j f_j^2)$.
\end{proof}

On the repository's reference dataset (200 reviewers, 10 items, a 60/40 split between two
camps; \texttt{sim/bridging\_irt\_dif.py}), the best of eight fits has
$\|\nabla L\|_\infty=7\times10^{-7}$. Every identity holds to that precision:
$\sum_j\iota_j=1.9\times10^{-8}$, $\sum_u\iota_u=4.2\times10^{-6}$,
$\lambda_b\sum_jf_j\iota_j=\lambda_f\sum_uf_u=-3.899\times10^{-3}$,
$\lambda_b\sum_uf_u\iota_u=\lambda_f\sum_jf_j=6.426\times10^{-2}$ and
$\sum_uf_u^2=\sum_jf_j^2=8.568$.

Community Notes also penalizes $\mu$~\cite{wojcik2022birdwatch}. Identity~\eqref{eq:zerosum}
then becomes $\sum_j\iota_j=\mu$ (with a common penalty weight), a constraint that is negligible
when hundreds of thousands of notes are fitted jointly. Isegoria leaves $\mu$ unpenalized and
fits one epoch's batch of tens to hundreds of items, so the identity binds.

\subsection{The gate is relative to the batch}

\begin{proposition}[Batch relativity]\label{prop:zerosum}\proved{}
Let every term of the minimum defining $B_j$ be a stationary point of \eqref{eq:loss}. Then
$\sum_j B_j\le 0$. Hence the mean bridge score of every batch is non-positive, and if
$\tau-\varepsilon>0$, every batch contains an item below the band: no batch passes the gate in
full.
\end{proposition}
\begin{proof}
Fix one fit $s_0$ among those in the minimum. Then $B_j\le\iota_j^{(s_0)}$ for every $j$, and
$\sum_j\iota_j^{(s_0)}=0$ by \cref{lem:gauge}. If every $B_j\ge\tau-\varepsilon>0$, then
$\sum_jB_j>0$, a contradiction.
\end{proof}

The proposition says that the gate does not test ``is this item approved across the divide?''.
It tests ``is this item approved across the divide at least $\tau$ more than the average
item fitted with it?''. \Cref{tab:relativity} shows the effect on the repository's dataset. We
score six consensus items twice: inside the original ten-item batch, which also contains a
true but divisive item, two partisan items and one mildly partisan item, and alone.

\begin{table}[t]
\centering\small
\caption{Bridge scores $B_j$ of the six consensus items of the repository simulation, fitted
inside the original ten-item batch and fitted alone (same ratings, same estimator). Threshold
$\tau=0.08$. Source: \texttt{scripts/levelA\_relativity.py}.}
\label{tab:relativity}
\begin{tabular}{@{}lrrrrrr@{}}
\toprule
Item & 01 & 02 & 04 & 05 & 06 & 07\\
\midrule
In the ten-item batch & $+0.104$ & $+0.081$ & $+0.085$ & $+0.101$ & $+0.077$ & $+0.082$\\
Alone                 & $+0.014$ & $-0.017$ & $-0.001$ & $+0.011$ & $-0.022$ & $-0.014$\\
\midrule
Passes in batch / alone & yes/no & yes/no & yes/no & yes/no & no/no & yes/no\\
\bottomrule
\end{tabular}
\end{table}

\begin{finding}\label{find:relativity}\empirical{}
Five of the six consensus items pass when the batch contains divisive items and none passes
when they are scored together without them. Their approval did not change; the reference
point did.
\end{finding}

Three consequences follow. First, pass rates depend on the composition of the epoch. A
threshold ``calibrated on the pilot'', as the design prescribes, is calibrated to the pilot's
mix of items. Second, the design documents found that the reference value of Community Notes
(0.40) was unusable and that $\tau\approx0.08$ worked, and that scores ``cluster near the
threshold''. Both observations are consistent with batch-centred scores. Third, whoever
influences the composition of a batch influences who passes. For example, weak ``decoy''
items lower the batch average and raise the relative score of a target. The admission lottery
and the proposal rate limit bound this lever but do not remove it.

Three corrections are possible, and we have not evaluated any of them \openq. One is to penalize $\mu$
and fit many epochs jointly, which dilutes the constraint. Another is to anchor the scale with reference items of known standing and
express $\tau$ relative to them; the golden items that already enter every review queue
(\cref{sec:levelc}) could serve. The third is to declare the gate relative on purpose, as an
explicit quota. The second fits the rest of the design best.

\subsection{The origin of the axis decides how majoritarian the score is}

The translation (G3) moves the origin of the reviewer axis and trades item intercept against
loading. This matters because the position of the origin decides what ``approval not explained
by the axis'' means. We make this precise in a two-camp model.

Consider camps $A$ and $B$ with shares $\pi_A,\pi_B$ and imbalance
$\kappa=\pi_B-\pi_A\ge0$ (so $B$ is the majority). For item $j$ let $m_{Aj},m_{Bj}$ be the
camps' mean ratings, $\mbar_j=\tfrac12(m_{Aj}+m_{Bj})$ the \emph{camp-balanced} mean and
$\Delta_j=\tfrac12(m_{Bj}-m_{Aj})$ the item's lean. The \emph{camp-size-weighted} mean, which
is what a plain average of all ratings estimates, is $\pi_Am_{Aj}+\pi_Bm_{Bj}=\mbar_j+\kappa\Delta_j$.
We say a stationary point has the \emph{two-camp form} if:
\begin{enumerate}[label=(T\arabic*)]
  \item every reviewer of camp $s$ has $f_u=\varphi_s$ and $\iota_u=0$, with half-distance
    $h=\tfrac12(\varphi_B-\varphi_A)>0$ and midpoint $c=\tfrac12(\varphi_A+\varphi_B)$; and
  \item the fit reproduces the camp means: $\mu+\iota_j+\varphi_s f_j=m_{sj}$ for $s\in\{A,B\}$ and every $j$.
\end{enumerate}

\begin{proposition}[Majoritarian leak]\label{prop:leak}\proved{}
Let a stationary point of \eqref{eq:loss} have the two-camp form. Let $S=\sum_j\Delta_j^2$ and
$S_{\Delta m}=\sum_j\Delta_j(\mbar_j-\mu)$. Then:
\begin{enumerate}[label=(\roman*)]
  \item $f_j=\Delta_j/h$ and $\iota_j=\mbar_j-\mu-t\,\Delta_j$, where $t=c/h$ is the offset of
    the origin from the midpoint between the camps;
  \item $t=\dfrac{\lambda_b S_{\Delta m}-\lambda_f\,n h^2\kappa}{\lambda_b S+\lambda_f\,n h^2}$;
  \item if $S_{\Delta m}=0$, then
  \begin{equation}\label{eq:leak}
    \mu+\iota_j=(1-\ell)\,\mbar_j+\ell\,(\pi_Am_{Aj}+\pi_Bm_{Bj}),\qquad
    \ell=\frac{1}{1+\rho},\qquad \rho=\frac{\lambda_b S}{\lambda_f\,n\,h^2};
  \end{equation}
  \item and, in that case, $\rho=\dfrac{\lambda_b}{\lambda_f}\,\gamma\sqrt{S/n}$ for some
    $\gamma\in\bigl[\sqrt{1-\kappa^2},\,1\bigr]$.
\end{enumerate}
\end{proposition}
\begin{proof}
Part (i) follows from subtracting and averaging the two equations of (T2). For (ii),
substitute (i) and $\sum_uf_u=n_A\varphi_A+n_B\varphi_B=nh(t+\kappa)$ into
\eqref{eq:gauge-rev}, which gives $\lambda_b(S_{\Delta m}-tS)/h=\lambda_f\,nh\,(t+\kappa)$, and
solve for $t$. Part (iii) is (i) with $t=-\kappa/(1+\rho)$ from (ii). For (iv), \eqref{eq:scale}
gives $nh^2(1+t^2+2\kappa t)=S/h^2$. On $t\in[-\kappa,0]$ the factor $1+t^2+2\kappa t$ lies in
$[1-\kappa^2,1]$, so $h^2=\sqrt{S/(n(1+t^2+2\kappa t))}$, and substituting into the definition
of $\rho$ gives (iv).
\end{proof}

The \emph{leak} $\ell$ is the fraction of the majority's weight that passes into the bridge
score. At $\ell=0$ the score is camp-balanced: each camp counts once, whatever its size. This is
the property that motivates bridging. At $\ell=1$ the score is the size-weighted mean, which
is exactly what a plain average gives. The asymmetric penalty pushes $\ell$ toward $0$, but
$\lambda_b\gg\lambda_f$ is not the relevant condition. The relevant quantity is
$\rho\approx(\lambda_b/\lambda_f)\sqrt{S/n}$, which compares how stiffly the penalty holds the
item intercepts (summed over the batch's leans) with how stiffly it holds the reviewer positions
(summed over the reviewers). At a fixed ratio $\lambda_b/\lambda_f$, the leak grows with the
number of reviewers and falls only with the square root of the batch's total lean.

Assumptions (T1)--(T2) are idealizations. With regularization no fit is exact, and reviewers of a camp
do not share one position. \Cref{tab:leak} therefore tests \eqref{eq:leak} on the full model:
all parameters free, rating noise and reviewer severity as in the repository simulation,
within-camp spread of positions, and nine ratings per reviewer. The design uses mirror pairs:
ten consensus items and five pairs of partisan items with identical quality and opposite lean
($\Delta=\pm0.36$), camps of 60\% and 40\% ($\kappa=0.2$). The measured leak is
$(\iota_{P^+}-\iota_{P^-})/(2\kappa\Delta)$, which is $0$ for a camp-balanced score and $1$ for a
size-weighted one.

\begin{table}[t]
\centering\small
\caption{Majoritarian leak $\ell$ of the bridge score: measured on the full model (mean of three
datasets; standard deviation in parentheses) against the prediction $1/(1+\rho)$ of
\cref{prop:leak} with $\gamma=1$. Mirror-pair design, 60/40 camps. The default penalties are
$\lambda_b/\lambda_f=5$. Source: \texttt{scripts/levelA\_leak.py}.}
\label{tab:leak}
\pgfplotstabletypeset[
  col sep=comma,
  columns={ratio,N,leak_mean,leak_sd,predicted},
  columns/ratio/.style={column name={$\lambda_b/\lambda_f$},int detect},
  columns/N/.style={column name={reviewers $n$},int detect},
  columns/leak_mean/.style={column name={measured $\ell$},fixed,fixed zerofill,precision=2},
  columns/leak_sd/.style={column name={(sd)},fixed,fixed zerofill,precision=2,
    postproc cell content/.append style={/pgfplots/table/@cell content/.add={(}{)}}},
  columns/predicted/.style={column name={predicted $1/(1+\rho)$},fixed,fixed zerofill,precision=2},
  every head row/.style={before row=\toprule,after row=\midrule},
  every last row/.style={after row=\bottomrule},
  every row no 6/.style={after row=\midrule},
  every row no 13/.style={after row=\midrule},
]{data/levelA_leak.csv}
\end{table}

\begin{finding}\label{find:leak}\empirical{}
At the default penalties most of the camp-size effect survives in the bridge score: the leak is
0.53 with 50 reviewers and rises to 0.87 with 3{,}200. The closed form predicts every configuration within
0.1 and overestimates the leak in most of them; within-camp spread, absent from (T1), adds to
$\sum_uf_u^2$ and hence to $\rho$. Raising $\lambda_b/\lambda_f$ to 500 keeps the leak below 0.1
up to 3{,}200 reviewers.
\end{finding}

\begin{figure}[t]
\centering
\begin{tikzpicture}
\begin{semilogxaxis}[
  width=0.58\linewidth, height=5.2cm,
  xlabel={reviewers $n$}, ylabel={leak $\ell$},
  ymin=-0.05, ymax=1.0, xmin=40, xmax=4000,
  legend style={font=\scriptsize, at={(1.03,1)}, anchor=north west, draw=none, fill=none},
  grid=major, grid style={gray!20}, tick label style={font=\scriptsize}, label style={font=\small}]
\addplot[only marks, mark=*, mark size=1.6pt, color=blue!70!black, skip coords between index={7}{21}]
  table[col sep=comma, x=N, y=leak_mean] {data/levelA_leak.csv};
\addlegendentry{$\lambda_b/\lambda_f=5$}
\addplot[only marks, mark=square*, mark size=1.5pt, color=orange!80!black, skip coords between index={0}{7}, skip coords between index={14}{21}]
  table[col sep=comma, x=N, y=leak_mean] {data/levelA_leak.csv};
\addlegendentry{$\lambda_b/\lambda_f=50$}
\addplot[only marks, mark=triangle*, mark size=1.8pt, color=green!45!black, skip coords between index={0}{14}]
  table[col sep=comma, x=N, y=leak_mean] {data/levelA_leak.csv};
\addlegendentry{$\lambda_b/\lambda_f=500$}
\addplot[thick, color=blue!70!black, forget plot, skip coords between index={7}{21}]
  table[col sep=comma, x=N, y=predicted] {data/levelA_leak.csv};
\addplot[thick, color=orange!80!black, forget plot, skip coords between index={0}{7}, skip coords between index={14}{21}]
  table[col sep=comma, x=N, y=predicted] {data/levelA_leak.csv};
\addplot[thick, color=green!45!black, skip coords between index={0}{14}]
  table[col sep=comma, x=N, y=predicted] {data/levelA_leak.csv};
\addlegendentry{prediction $1/(1+\rho)$ (lines)}
\end{semilogxaxis}
\end{tikzpicture}
\caption{Majoritarian leak against the number of reviewers (points: full model; lines:
$1/(1+\rho)$ from \cref{prop:leak}). Data as in \cref{tab:leak}.}
\label{fig:leak}
\end{figure}

\Cref{prop:leak} also explains an asymmetry in the repository simulation. There $n=200$,
$S\approx0.39$ and $\kappa=0.2$, so $\rho\approx0.22$ and $\ell\approx0.82$. The majority's
partisan item and the minority's mirror-image item have the same quality and opposite lean. In
the full fit they score $-0.155$ and $-0.260$. The proposition predicts a gap of
$\ell\kappa(\Delta_{08}-\Delta_{09})\approx0.82\times0.2\times0.72\approx0.12$; the measured
gap is $0.105$. Both items fail here. For items closer to $\tau$, however, which of two mirror
images passes depends on which camp is larger, which is what bridging is meant to prevent.

The normative content of bridging, ``each side counts once'', therefore sits in a gauge. The
objective reaches it only in the limit $\rho\to\infty$, and with more than two groups or a
continuum of positions the ``midpoint'' is itself a choice. We see three remedies
\openq:
\begin{enumerate}[label=(\alph*)]
  \item scale the penalty ratio with $\sqrt{n/S}$; for example, $\rho\ge9$ keeps $\ell\le0.1$;
  \item fix the gauge explicitly, for instance by imposing $\sum_jf_j\iota_j=0$, the
    $\rho\to\infty$ limit of \eqref{eq:gauge-rev}. The origin then depends on $S_{\Delta m}$, the
    correlation between lean and approval within the batch, which is a new lever for anyone
    who controls the batch's composition;
  \item replace the intercept by a gauge-invariant statistic of the predictions, such as the
    smaller of the mean predicted ratings $\hat r_{uj}$ among reviewers in the lowest and in the
    highest decile of $f_u$. Predictions are invariant under (G1)--(G5), so any function of
    them is too, and a minimum across the spectrum does not depend on camp sizes.
\end{enumerate}

\subsection{What ratings cannot distinguish}

\begin{proposition}[Observational equivalence]\label{prop:equivalence}\proved{}
Let $D$ be any decision rule, deterministic or randomized, that depends on the data only
through the assignment $\Omega$ and the ratings $(r_{uj})$. If two items induce the same
distribution of ratings from every reviewer, then $D$ has the same distribution of outcomes on
both.
\end{proposition}
\begin{proof}
The distribution of $D$'s output is a function of the joint distribution of its inputs, which is
the same for both items.
\end{proof}

The proposition is elementary, but it settles the most serious limitation of Level~A. A true
fact that one camp rejects and a false claim that the other camp accepts can polarize the same
camps in the same way. When they do, no function of the ratings separates them, and every
threshold trades lost true items for admitted partisan ones. The repository simulation shows
this: at every $\tau$ in $[0.02,0.12]$ the true but divisive item fails, while no partisan item
passes. The remedy has to bring in information the ratings do not contain. That is the role of
the appeal to evidence, which sends a polarization rejection to the pilots at the author's risk.
\Cref{sec:dif-not-bias} discusses how far that remedy reaches.

\subsection{Capture cost in the reference simulation}

In the repository simulation, 40 members of the majority camp rate the majority's partisan
item at $1$, and members of the opposing camp are added one group at a time. The bridge score
moves from $-0.14$ (no opposing raters) through $-0.10$, $-0.06$, $-0.02$, $+0.02$ and $+0.07$ to
$+0.13$ (10, 20, 30, 40, 55 and 70 opposing raters). The crossing of $\tau=0.08$ therefore lies
between 55 and 70 of the 80 opposing reviewers (69--88\%); linear interpolation gives about 57
(71\%). The design documents quote ``70 of 80 (87\%)'', which is the upper end of this bracket
\empirical. This is one dataset and one seed. A capture-cost curve with confidence bands is
part of the evaluation plan (\cref{sec:plan}). \Cref{prop:leak} adds that votes from the
majority's own camp are not neutralized: they act through the leak.

\subsection{Further identifiability conditions}

The sign of the axis is fixed as described above. If the bipartite graph of $\Omega$ is
disconnected, each component has its own gauge and the scores of different components are not
comparable. Random assignment with $k\ge7$ reviewers per item makes this unlikely, but the
engine does not check it. The specification states that reviewers with fewer than $n_{\min}=30$
ratings should not define the axis. This rule and the two-dimensional model are not
implemented; the second is deferred by design decision D31~\cite{isegoria-repo}.
